\chapter{Scenarios}
\label{ch:scenarios}

A scenario is the primary executable narrative unit in LeanFM. It is a
named, bounded stream of timestamped events that together describe one
interaction or one execution trace of interest (for example: a successful
quote request, an OAuth exchange, or a two-player 9×9 game). Scenarios are the
basis for generated interaction diagrams, state machines, and the numeric
and histogram charts used to justify design decisions.

\section{Canonical event fields}
Every scenario event is an envelope with a small set of durable fields used by
the tooling and by requirements. Typical fields include:

\begin{itemize}
  \item id: stable event identifier (unique within the project).
  \item prior\_ids: zero or more predecessor event ids that establish the
    partial order for the scenario (forks and joins are permitted).
  \item cause\_id: the direct causal parent for paired exchanges (for example,
    an HTTP response names the request event that caused it).
  \item session, scenario: the pair that bounds the interaction and groups
    events for derived views.
  \item src, dst: the actor instance names for sender and receiver lifelines.
  \item timeAt: ISO-8601 timestamp for wall-clock ordering and latency
    measurement.
  \item message: the observable message or terminal label for the event.
\end{itemize}

These fields support the two distinct orderings used throughout LeanFM: \emph{prior}-based partial order (for concurrency and commutation) and the
separate \emph{cause} pointer used to pair request/response-style exchanges.
Do not conflate the two: cause\_id pairs logical call/response, while prior\_ids
control allowed interleavings and render order in interaction diagrams.

\section{Derived outputs and provenance}
By default the book and server generate the following per `(session,scenario)`:
\begin{itemize}
  \item an interaction diagram (lifelines and arrows rendered from src/dst/\(message\)),
  \item a residual state machine (observed states and message-labeled transitions),
  \item latency and throughput charts derived from timeAt and prior start/stop pairs,
  \item outcome distributions and other question-specific reductions.
\end{itemize}

Every generated view carries provenance metadata: the originating prompt or
requirement id, the reducer used to compute the view (for example: partition by
(session,scenario) and render an order graph from prior pointers), the unit of
measurement, and the scenario boundary. When prose describes an interaction
that a generated view can render, replace most of that prose with the view and
a short interpretation during the book cleanup phase (see SKILL.md).

\section{Example: minimal two-player game scenario}
Below is an illustrative minimal fragment of a scenario event list for a
9×9 game (timestamps and ids are abbreviated):
\begin{verbatim}
{ "id":"evt-0010", "prior_ids":["evt-0009"], "cause_id":null,
  "timeAt":"2026-10-17T12:00:30Z", "session":"s-g1", "scenario":"g1",
  "src":"alice","dst":"server","message":"Move E5 (B)" }
{ "id":"evt-0011", "prior_ids":["evt-0010"], "cause_id":"evt-0010",
  "timeAt":"2026-10-17T12:00:30Z", "session":"s-g1", "scenario":"g1",
  "src":"server","dst":"alice","message":"MoveAccepted E5" }
\end{verbatim}

This fragment is sufficient for an interaction diagram and for latency
measurement between a client's Move and the server's MoveAccepted reply.

\section{Guidance for authors}
When writing or editing the book:
\begin{itemize}
  \item prefer concrete scenario examples that the generator can render; include
    tiny, runnable traces rather than only prose descriptions.
  \item record the intended derived outputs (interaction diagram, state
    machine, charts) alongside the scenario so the reader can check the
    generated views against the prose claim.
  \item avoid mixing implementation transport details into Requirements. Use
    Implementation.lean for transport and runtime mapping after requirements
    acceptance.
\end{itemize}

Scenario prose and small example traces greatly increase the book's value for
codex and human reviewers. This chapter is intended as a compact reference and
as a reminder to bind generated views to the prompt and reducer that justified
them.